\documentclass[10pt,a4paper]{amsart}
\usepackage[margin=19mm]{geometry}
\usepackage{amsmath,amssymb,mathtools}
\usepackage[scr=rsfs,cal=euler]{mathalfa}
\usepackage{xfrac}
\usepackage[unicode,hidelinks,bookmarks=false]{hyperref}
\newcommand{\ie}{i.e.\ }
\newcommand{\cf}{cf.\ }
\newcommand{\resp}{respectively\ }
\newcommand{\Subsection}{\subsection}
\providecommand{\nobreakdash}{\nobreak\hbox{-}}
\setlength{\emergencystretch}{2em}
\multlinegap0pt
\usepackage{bm}
\usepackage{bbm}
\usepackage{dsfont}
\usepackage{xspace}
\usepackage{cancel}
\usepackage{mathtools}
\usepackage{amsthm}
\makeatletter
\@ifundefined{theorem}{\newtheorem{theorem}{Theorem}}{}
\makeatother
\title[Cohen, Levitzki, Hilbert Basis, and Lasker-Noether Theorems ]{Cohen, Levitzki, Hilbert Basis, and Lasker-Noether Theorems for Nil-S-Noetherian Rings}
\author{Aman Pandey, Ajim Uddin Ansari,, Tushar Singh}
\date{Source: arXiv:2606.27024v1 (CC BY 4.0)}
\begin{document}
\maketitle

\noindent\emph{Source and licence.} Cohen, Levitzki, Hilbert Basis, and Lasker-Noether Theorems for Nil-S-Noetherian Rings. Version arXiv:2606.27024v1 (CC BY 4.0), distributed under the Creative Commons Attribution 4.0 International licence, as recorded in the arXiv licence metadata for this exact version and shown on the abstract page https://arxiv.org/abs/2606.27024v1. Full rights notice: licenses/arxiv-2606.27024-CC-BY-4.0.txt.\par\smallskip

\noindent\textit{Editorial scope.} Counted unit: the Theorem whose source label is \texttt{Hilbert-2} in the article (source0.tex lines 330--332, source label \texttt{Hilbert-2}), delivered together with the article's own complete original proof of it (source0.tex lines 334--362, one \texttt{proof} environment of 29 lines). No mathematics is written, repaired or re-derived: statement and proof are the article's own text line for line. Required context retained uncounted: cohen (lines 245--251). Every definition, display and label of those lines is retained unchanged. Omitted from this package and not redistributed: 1--244, 252--329, 333--333, 363--467. Nothing inside a retained excerpt is omitted or altered: every retained body is delivered exactly as the article's lines are written, so no omission receipt of any kind accompanies this package. Citations: the retained excerpts carry no citation command at all, so no bibliography entry of the article is transcribed and no reference is redistributed. Reference adaptation: none was needed and none was made; every reference inside the delivered unit resolves inside it, so no unresolved reference or citation is produced. Printed slips kept literal: no printed word, symbol, hypothesis or number of the counted statement or of its proof was altered. Layout and class: the article's own class is \texttt{article} with packages including amsmath, amsthm, amssymb, latexsym, fontenc, authblk, lipsum, hyperref, lineno, upgreek, tikz; the wrapper is the lane's pinned \texttt{amsart} editorial wrapper at 10pt on A4, which restates the theorem-like environments the retained text needs and adds the article's own macro definitions that the retained text needs (none), with extra wrapper packages (bm, bbm, dsfont, xspace, cancel, mathtools). The delivered numbering is the unit's own. Assets: no retained span contains a figure environment, a graphics-inclusion command, a PDF-page inclusion, a table or a TikZ picture; no image file is copied into this package, the package declares no asset dependency and no image file is redistributed with it. Licence: version arXiv:2606.27024v1 of this article is distributed under the Creative Commons Attribution 4.0 International licence, as recorded in the arXiv licence metadata for this exact version and shown on the abstract page https://arxiv.org/abs/2606.27024v1. A full rights notice is delivered at licenses/arxiv-2606.27024-CC-BY-4.0.txt. Characters: every retained excerpt is pure ASCII, so the delivered engine, XeLaTeX with its default Latin Modern text font, reports no missing character.
\begin{theorem} \label{cohen}\textbf{(Cohen theorem for Nil-$S$-Noetherian rings).}
Let $R$ be a ring. Then the following statements are equivalent:
\begin{enumerate}
    \item $R$ is Nil-$S$-Noetherian. 
    \item Every prime ideal $P \subseteq Nil(R)$ with $P \cap S = \emptyset$ is $S$-finite.
\end{enumerate}
\end{theorem}

\begin{theorem} \label{Hilbert-2}
	Let $R$ be a ring and $S \subseteq R$ be an anti-Archimedean multiplicative set  consisting of non-zero divisors. If $R$ is Nil-$S$-Noetherian, then the power series ring $R[[x]]$ is also Nil-$S$-Noetherian.
\end{theorem}

\begin{proof}
Let $P$ be a nil prime ideal of $R[[x]]$ such that $P\cap S\neq\emptyset$.	 Let
	$ P(0) = \{f(0) \mid f(x) \in P\} $. 
	Let $f= \sum_{i=0}^{\infty}a_{i}x^{i} \in P$. Then $f$ is nilpotent as $P$ is nil ideal. This implies that each coefficient $a_{n}$ is nilpotent, in particular, its constant term $a_{0} = f(0)$ is nilpotent.
	Consequently, $P(0)$ is a nil ideal of $R$. This implies that $P(0)$ is $S$-finite since $R$ is Nil-$S$-Noetherian. Thus, there exists $s \in S$ and a finite set of polynomials $\{p_1(x), p_2(x), \dots, p_n(x)\} \subseteq P$ such that
	$ sP(0) \subseteq \langle p_1(0), p_2(0), \dots, p_n(0) \rangle \subseteq P(0) $.
	If $x \in P$, then we can write $P = P(0) + xR[[x]]$. It follows that
	 $sP \subseteq \langle p_1(0), p_2(0), \dots, p_n(0), x \rangle \subseteq P$.
	Hence $P$ is $S$-finite.
	 Now suppose  $x \notin P$.
	Let $f \in P$. Then $f(0) \in P(0)$. Write
	$ sf(0) = \sum_{i=1}^{n} r_{0i} p_i(0) \quad \text{for some } r_{0i} \in R$.
	This implies that
	$ sf -  sf(0) = sf- \sum_{i=1}^{n} r_{0i} p_i(0) = xf_1 $
	for some $f_1 \in R[[x]]$. Since $xf_1 \in P$ and $x \notin P$ (as $P$ is a prime ideal), it must be that $f_1 \in P$. 	
	By repeating this iterative process, we can similarly write
	$sf_1 - \sum_{i=1}^{n} r_{1i} p_i(0) = xf_2 $,
	where $f_2 \in P$. Continuing this construction, we obtain the representation
	$ f = \sum_{i=1}^{n} p_i(0) \left( \sum_{j=0}^{\infty} \frac{r_{ji}}{s^{j+1}} x^j \right) $.
	Note that the coefficients $\frac{r_{ji}}{s^{j+1}}$ belong to the localization $S^{-1}R$. Since $S$ is given to be an anti-Archimedean set, we have
	$ \left( \bigcap_{i=1}^{\infty} s^i R \right) \cap S \neq \emptyset $.
	Let $t \in \left( \bigcap_{i=1}^{\infty} s^i R \right) \cap S$. Then for each $j$, $t = s^j c_j$ for some $c_j \in R$. 
	Multiplying the expression for $f$ by $t$, we get
	$ tf = \sum_{i=1}^{n} p_i(0) \left( \sum_{j=0}^{\infty} c_{j+1} r_{ji} x^j \right) \in \langle p_1(0), \dots, p_n(0) \rangle $.
	This implies that
	$ tP \subseteq \langle p_1(0), p_2(0), \dots, p_n(0) \rangle \subseteq P$.
	Thus $P$ is $S$-finite. Hence by Theorem~\ref{cohen}, $R[[x]]$ is a Nil-$S$-Noetherian ring.
	
\end{proof}

\end{document}
